\begin{thebibliography}{99}
	
	\section*{Wydawnictwa zwarte}
	
	\bibitem[1]{1}
	I.~Sommerville,  
	\emph{Inżynieria oprogramowania}, wyd.~10,  
	Wydawnictwo Naukowe PWN, Warszawa, 2016.
	
	\bibitem[2]{2}
	R.C.~Martin,  
	\emph{Czysta architektura. Struktura i design oprogramowania. Przewodnik dla profesjonalistów},  
	Wydawnictwo Helion, Gliwice, 2018.
	
	\section*{Netografia}
	
	\bibitem[www1]{www1}
	ITU-T,  
	\emph{Recommendation ITU-T Y.2060 (06/2012): Overview of the Internet of things},  
	International Telecommunication Union, Genewa, 2012,  
	\url{https://www.itu.int/rec/T-REC-Y.2060-201206-I}  
	[dostęp: 08.11.2026].
	\bibitem[www2]{www2}
	S.~Sinha,  
	\emph{State of IoT 2025: Number of connected IoT devices growing 14\% to 21.1 billion globally},  
	IoT Analytics, 2025,  
	\url{https://iot-analytics.com/number-connected-iot-devices/}  
	[dostęp: 08.11.2026].
	
	\bibitem[www3]{www3}
	Atlassian,  
	\emph{User stories},  
	\url{https://www.atlassian.com/pl/agile/project-management/user-stories}  
	[dostęp: 25.05.2026].
	
	\bibitem[www4]{www4}
	Perforce,  
	\emph{What Are Non-Functional Requirements?},  
	\url{https://www.perforce.com/blog/alm/what-are-non-functional-requirements-examples}  
	[dostęp: 25.05.2026].
	
\end{thebibliography}